% BHU PYQ -- Professional Exam Template
\documentclass[11pt,a4paper]{article}

\usepackage[a4paper,top=2.5cm,bottom=2.5cm,left=2.8cm,right=2.8cm,headheight=14pt]{geometry}
\usepackage{amsmath,amssymb}
\usepackage[T1]{fontenc}
\usepackage[utf8]{inputenc}
\usepackage{lmodern}
\usepackage{microtype}
\usepackage{fancyhdr}
\usepackage{enumitem}
\usepackage{listings}
\usepackage{hyperref}
\usepackage{lastpage}
\usepackage{parskip}

\hypersetup{colorlinks=true,urlcolor=black,linkcolor=black,
  pdftitle={BVCA-101: Introduction To Computers -- BHU PYQ}}

\pagestyle{fancy}
\fancyhf{}
\renewcommand{\headrulewidth}{0.4pt}
\renewcommand{\footrulewidth}{0.4pt}
\fancyhead[L]{\small\textit{Banaras Hindu University}}
\fancyhead[R]{\small\textit{BVCA-101: Introduction To Computers}}
\fancyfoot[C]{\small Page \thepage\ of \pageref{LastPage}}

\lstset{
  basicstyle=\small\ttfamily,
  frame=single,
  breaklines=true,
  columns=flexible,
  keepspaces=true
}

\newlist{parts}{enumerate}{2}
\setlist[parts,1]{label=(\alph*),leftmargin=*,itemsep=4pt,topsep=3pt}
\setlist[parts,2]{label=(\roman*),leftmargin=*,itemsep=2pt,topsep=2pt}
\newcommand{\pts}[1]{\hfill\textit{[#1]}}

\begin{document}

\begin{center}
  {\large\bfseries BANARAS HINDU UNIVERSITY}\\[2pt]
  {\normalsize B.Voc. Semester I Examination, 2022-23}\\[6pt]
  \rule{0.8\linewidth}{0.5pt}\\[6pt]
  {\large\bfseries Paper: BVCA-101 --- Introduction To Computers}\\[4pt]
  {\normalsize Time: Three Hours \hfill Maximum Marks: 70}
\end{center}

\rule{\linewidth}{0.4pt}

\smallskip
\noindent\textit{Instructions: Attempt all sections. Questions of Sections A, B and C carry 15, 5 and 1 marks each respectively. Symbols have their usual meanings.}

\bigskip

\bigskip\noindent{\large\bfseries SECTION --- A}

\rule{\linewidth}{0.4pt}\smallskip

\noindent\textbf{Long Answer Questions} (Attempt all questions)\pts{2 $\times$ 15 = 30}
\begin{enumerate}[label=\arabic*.,leftmargin=*,itemsep=12pt]
  \item Discuss the generations of computers based on their characteristics and processing speed.
  \begin{center}\textbf{OR}\end{center}
  Based on the block diagram of a computer, discuss its different hardware components.
  \item What is a Compiler? Describe how a compiler differs from an interpreter in detail.
  \begin{center}\textbf{OR}\end{center}
  Define operating system and give a detailed account on different types of operating systems.
\end{enumerate}

\bigskip\noindent{\large\bfseries SECTION --- B}

\rule{\linewidth}{0.4pt}\smallskip

\noindent\textbf{Short Answer Questions} (Attempt any six questions)\pts{6 $\times$ 5 = 30}
\begin{enumerate}[label=\alph*.,leftmargin=*,itemsep=8pt]
  \item What is DOS? Describe the internal and external commands of DOS.
  \item Provide the differences between procedural and object-oriented programming paradigms.
  \item Write short notes on the following:
  \begin{parts}
    \item Overloading
    \item Polymorphism
  \end{parts}
  \item Define the mail merge feature in MS-Office Word.
  \item What are the advantages and disadvantages of a Wide Area Network (WAN)?
  \item What are the differences between primary memory and secondary memory?
  \item Explain the different types of common network topologies.
  \item What is formatting in Excel? Explain the formatting features available in Excel for numeric data.
\end{enumerate}

\bigskip\noindent{\large\bfseries SECTION --- C}

\rule{\linewidth}{0.4pt}\smallskip

\noindent\textbf{Multiple Choice Questions} (Answer all questions)\pts{10 $\times$ 1 = 10}
\begin{enumerate}[label=\arabic*.,leftmargin=*,itemsep=8pt]
  \item Who is the father of Computers?
  \begin{parts}
    \item James Gosling
    \item Charles Babbage
    \item Dennis Ritchie
    \item Bjarne Stroustrup
  \end{parts}
  \item Which of the following is the correct abbreviation of COMPUTER?
  \begin{parts}
    \item Commonly Occupied Machines Used in Technical and Educational Research
    \item Commonly Operated Machines Used in Technical and Environmental Research
    \item Commonly Oriented Machines Used in Technical and Educational Research
    \item Commonly Operated Machines Used in Technical and Educational Research
  \end{parts}
  \item What is the full form of CPU?
  \begin{parts}
    \item Computer Processing Unit
    \item Computer Principle Unit
    \item Central Processing Unit
    \item Control Processing Unit
  \end{parts}
  \item Which of the following computer languages is written in binary codes only?
  \begin{parts}
    \item Pascal
    \item Machine language
    \item C
    \item C\#
  \end{parts}
  \item Which of the following is the smallest unit of data in a computer?
  \begin{parts}
    \item Bit
    \item KB
    \item Nibble
    \item Byte
  \end{parts}
  \item Which of the following services allows a user to log in to another computer somewhere on the Internet?
  \begin{parts}
    \item e-mail
    \item UseNet
    \item Telnet
    \item FTP
  \end{parts}
  \item The software designed to perform a specific task:
  \begin{parts}
    \item Synchronous Software
    \item Package Software
    \item Application Software
    \item System Software
  \end{parts}
  \item What do you call a specific instruction designed to do a task?
  \begin{parts}
    \item Command
    \item Process
    \item Task
    \item Instruction
  \end{parts}
  \item Public domain software is usually:
  \begin{parts}
    \item System supported
    \item Source supported
    \item Community supported
    \item Programmer supported
  \end{parts}
  \item OSS stands for:
  \begin{parts}
    \item Open System Service
    \item Open Source Software
    \item Open System Software
    \item Open Synchronized Software
  \end{parts}
\end{enumerate}

\vfill
\rule{\linewidth}{0.4pt}
\begin{center}\small
\href{https://bhu-exam.vercel.app/}{Click here to visit our website}\\[4pt]\href{https://me-aryan.vercel.app/}{Click here to visit our contributor}\\[4pt]\href{https://quashan-me.vercel.app/}{Click here to visit our contributor}
\end{center}
\end{document}
